\documentclass[conference]{IEEEtran}
\usepackage{cite}
\usepackage{amsmath,amssymb,amsfonts}
\usepackage{algorithmic}
\usepackage{graphicx}
\usepackage{textcomp}
\usepackage{xcolor}
\usepackage{booktabs}
\usepackage{multirow}
\usepackage{float}
\usepackage{url}

\begin{document}

\title{Codex (GPT-5.4) as an Agentic AI Coding Assistant:\\A Practitioner Benchmark Study in Production Software Development}

\author{\IEEEauthorblockN{Manee \& Sushma}
\IEEEauthorblockA{\textit{ACME Development Team}}}

\maketitle

\begin{abstract}
Codex (OpenAI, 2026), powered by the unified GPT-5.4 architecture, is OpenAI's agentic coding assistant featuring cloud-native sandboxed execution, git worktree isolation, purpose-trained context compaction, and a dual-layered security model. We evaluate it as an AI coding assistant through a practitioner benchmark involving real feature implementation in a production SaaS codebase, assessing eight dimensions: workflow integration, code generation quality, contextual reasoning, architectural adherence, problem-solving capability, documentation generation, operational characteristics, and overall productivity impact.

Using a 26-criterion rubric (weighted mean: 4.50/5.00, 90.0\%), published benchmark triangulation verified against primary sources, and practitioner community sentiment analysis, we find that Codex excels at analytical problem solving, fast-mode execution velocity, sustained long-horizon task completion, and technical debt remediation. Its SWE-Bench Pro leadership (57.7\% vs.\ Claude Opus 4.6's 53.4\%) reflects superior novel algorithmic reasoning, while its OSWorld-Verified score (75.0\%) makes it the first general-purpose AI model to definitively exceed the human expert threshold (72.4\%) in desktop automation. Key limitations include a documented tendency toward destructive execution under underspecified prompts, security and IP constraints inherent to cloud-hosted architectures, and rate-limiting friction on professional Pro-tier subscriptions. These findings offer practical guidance for teams evaluating agentic AI coding assistants for production use.
\end{abstract}

\begin{IEEEkeywords}
Codex, GPT-5.4, agentic AI, coding assistant, benchmark, git worktrees, cloud execution, context compaction, SWE-Bench Pro
\end{IEEEkeywords}

\section{Introduction}
AI-assisted software development has entered a qualitatively new phase in 2026. The industry has moved decisively beyond single-context autocomplete assistants toward highly autonomous, long-horizon agentic systems capable of orchestrating complex multi-file codebase refactoring, executing deterministic terminal operations, and navigating desktop environments natively. At the epicenter of this convergence is OpenAI's release of the unified GPT-5.4 architecture, which fully absorbed the capabilities of the previously dedicated GPT-5.3-Codex model into a single frontier system, paired with a substantially revamped Codex standalone application.

Codex distinguishes itself from competing agentic tools through several architectural innovations that have no direct analogues in locally-executed alternatives: cloud-native sandboxed container execution, git worktree isolation enabling true parallel task processing across multiple projects, and a purpose-trained context compaction mechanism that enables coherent multi-session operation far exceeding standard context window limits. These features collectively position Codex as the primary execution-focused counterpart to Claude Code in the frontier agentic tools landscape.

\subsection{Codex: Technical Overview}
Codex can be accessed through multiple channels:
\begin{itemize}
    \item \textbf{Standalone application}: A full-featured desktop app superseding the legacy CLI tool, supporting customizable skills, automation pipelines, and native Machine Context Protocol (MCP) server toggles \cite{openai_codex_docs}
    \item \textbf{VS Code extension}: Fully compatible with advanced forks including Cursor and Windsurf \cite{openai_codex_docs}
    \item \textbf{JetBrains IDE integration}: Native support for IntelliJ, PyCharm, Rider, and WebStorm \cite{openai_codex_docs}
    \item \textbf{CLI (open-source)}: Available on GitHub, configurable via layered \texttt{.codex/config.toml} hierarchies with predictable resolution precedence
\end{itemize}

The model's most distinctive architectural characteristics include:
\begin{itemize}
    \item \textbf{Cloud worktree execution}: Tasks run in isolated OpenAI-managed containers with per-task git worktree separation, enabling parallel execution across multiple projects simultaneously without consuming local developer machine resources \cite{openai_cloud_envs}
    \item \textbf{Native context compaction}: GPT-5.4 is purpose-trained to operate coherently across multiple context windows---not mere summarization, but a model-level architectural capability enabling sessions exceeding 24 hours with approximately 30\% fewer thinking tokens than predecessor models \cite{openai_compaction}
    \item \textbf{1.05M token context window}: The largest context window among frontier coding models, complemented by a Tool Search feature that dynamically loads tool definitions and API schemas on demand, reducing effective input token consumption by approximately 47\% for agentic tasks \cite{openai_compaction}
    \item \textbf{SWE-Bench Pro leadership}: 57.7\% on novel, non-gameable engineering tasks, outperforming Claude Opus 4.6 (53.4\%) by a 4.3 percentage point margin reflecting genuine superiority in first-principles algorithmic reasoning \cite{morphllm2026}
    \item \textbf{Dual-layered security model}: Sandbox Mode (technical execution boundaries) and Approval Policy Engine (human oversight checkpoints) operating as independent, composable controls with platform-native enforcement via macOS Seatbelt, Linux \texttt{bubblewrap}, and Windows WSL2 \cite{openai_codex_docs}
    \item \textbf{Fast Mode}: Delivers up to 1.5$\times$ faster token velocity using the identical underlying model intelligence, without quality trade-offs, critical for maintaining developer flow state on well-defined tasks \cite{holter2026}
    \item \textbf{AGENTS.md project instruction hierarchy}: A deterministic, version-controllable file-based instruction system that the GPT-5.4 model is explicitly fine-tuned to follow, providing more predictable constraint adherence than free-form session prompting
\end{itemize}

\subsection{Research Questions}
This paper addresses three primary research questions:

\textbf{Q1}: How effectively does Codex (GPT-5.4) function as an AI coding assistant in a structured, convention-driven production codebase, specifically in terms of analytical problem solving, sustained execution, and architectural reasoning?

\textbf{Q2}: What are the observable operational characteristics of Codex---including its cloud execution model, context compaction behavior, rate-limiting profile, and cost-effectiveness---that affect its suitability for daily professional use?

\textbf{Q3}: How do Codex's unique architectural features---git worktrees, cloud-native sandboxing, native compaction, and Tool Search---provide a competitive edge over single-context and locally-executed agentic alternatives?

\section{Background and Related Work}

\subsection{Benchmarks for AI Code Generation}
The evaluation of AI coding assistants has evolved through several benchmark generations. HumanEval \cite{chen2021evaluating} established an early standard through 164 Python programming problems but has been widely criticized for narrow scope and susceptibility to data contamination \cite{fryer2024hallucination}. SWE-bench \cite{jimenez2024swebench} advanced the field significantly by using real GitHub issues requiring multi-file patches.

The SWE-Bench ecosystem has itself bifurcated. SWE-Bench Verified, which tests standardized and heavily vetted pull request resolution, shows near-parity between GPT-5.4 ($\sim$80.0\%) and Claude Opus 4.6 (81.4\%), with Opus 4.6's advantage reflecting superior meticulous execution of well-defined patterns \cite{vellum2026}. The more consequential metric is \textbf{SWE-Bench Pro}, specifically engineered to resist training data memorization by forcing agents to rely on novel logical reasoning rather than pattern matching. On this benchmark, GPT-5.4 achieves 57.7\%, outperforming Claude Opus 4.6's 53.4\%---a 4.3 percentage point advantage that reflects genuine superiority in algorithmic synthesis and structural deduction when confronted with problems outside the training distribution \cite{morphllm2026}.

OSWorld-Verified extends evaluation into autonomous desktop navigation, measuring an AI agent's ability to interact with desktop GUIs, analyze high-resolution screenshots, and complete multi-step workflows. GPT-5.4's score of 75.0\% marks the first time a general-purpose AI has definitively exceeded the human expert baseline of 72.4\% on complex desktop automation tasks, powered by a vision architecture that processes images containing over 10 million pixels without destructive compression \cite{nxcode2026}.

\subsection{Developer Productivity Studies}
A landmark randomized controlled trial involving over 4,000 developers found that GitHub Copilot reduced mean task completion time by approximately 55\% for new coding tasks, with junior developers showing disproportionate gains \cite{peng2023copilot}. McKinsey Global Institute reported a 26\% productivity increase in developer output when using AI tools, measured as story points per sprint \cite{mckinsey2023}.

A contrarian finding from METR \cite{metr2025} observed a 19\% \textit{increase} in task completion time for experienced developers on complex, unfamiliar codebases. This suggests that productivity gains from AI coding tools are not uniform and depend heavily on codebase familiarity, task clarity, and the degree to which developers retain architectural ownership. These findings underscore the importance of the \texttt{AGENTS.md} discipline: well-specified constraints reduce the ambiguity that drives both METR-style slowdowns and Codex's documented destructive execution tendency.

\subsection{Comparative Landscape}
As of 2026, agentic coding tools have bifurcated along two key architectural philosophies:
\begin{itemize}
    \item \textbf{IDE-first vs.\ cloud-first}: Cursor and Copilot (deep editor integration, single-context) vs.\ Claude Code and Codex (repository-wide autonomy, with Codex uniquely offering cloud-native container execution)
    \item \textbf{Parallel orchestration vs.\ sustained sequential execution}: Claude Code's Agent Teams framework (5--6 parallel specialized agents with proactive QA dispatch) vs.\ Codex's sustained long-horizon sequential execution with delegation to cheaper submodels for routine subtasks
\end{itemize}

Codex occupies a distinctive position: the only frontier tool offering true cloud-native parallel task execution via git worktrees, where each task runs in an isolated OpenAI-managed container. This architecture has no direct analogue in Claude Code's local execution model and represents a fundamental operational difference for teams managing multiple concurrent workstreams.

\section{Methodology}

\subsection{Evaluation Design}
This study uses a practitioner observational design \cite{creswell2018mixed}: developers used Codex as the primary AI assistant over a structured evaluation period for real production tasks, and the resulting experience and outputs were evaluated using three independent instruments, supplemented by cross-verified external benchmark triangulation and practitioner community sentiment analysis drawn from developer forums, technical blogs, and comparative reviews.

\subsection{Evaluation Environment}
The evaluation was conducted on a production SaaS platform for social media content scheduling. The codebase is a monorepo comprising:
\begin{itemize}
    \item \textbf{Backend}: Python FastAPI following strict Clean Architecture layering (API $\rightarrow$ Controller $\rightarrow$ UseCase $\rightarrow$ Domain $\rightarrow$ Infrastructure)
    \item \textbf{Frontend}: Flutter mobile application with Clean Framework and Riverpod v3 state management
\end{itemize}

The codebase represents an authentic evaluation environment: non-trivial architectural conventions, strict layer boundaries, an established pattern library Codex was expected to discover and follow, and a running Docker Compose stack for integration verification. An \texttt{AGENTS.md} file at the repository root provided explicit architectural conventions, key file locations, build and test commands, and negative constraints.

\subsection{Evaluation Tasks}
Codex was evaluated across multiple task categories over a structured evaluation period:

\textbf{Phase 1}: Setup and initial configuration, Fast Mode vs.\ Standard Mode comparison, cloud worktree initialization and parallel task execution test

\textbf{Phase 2}: Context retention and \texttt{AGENTS.md} constraint fidelity testing, context compaction behavior monitoring over extended multi-hour sessions

\textbf{Phase 3}: Full-stack feature implementation (Pinterest default board selection, 24 files modified/created) via cloud sandbox execution with git worktree isolation

\textbf{Phase 4}: Sustained execution sprint evaluation---technical debt identification and automated legacy refactoring over multi-hour autonomous sessions

\textbf{Phase 5}: Bug identification in existing codebase, cross-model audit workflow (Codex reviewing Claude-generated code), and documentation generation assessment

\subsection{Evaluation Instruments}

\textbf{Instrument 1: Quantitative Rubric}. Developers rated Codex on a 26-criterion rubric spanning eight categories: Workflow and Setup, Code Generation Quality, Context and Memory Management, Architecture and Planning, Problem Solving, Operational Characteristics, Documentation, and Overall Productivity Impact. Each criterion was scored 1 (critically deficient) to 5 (excellent).

\textbf{Instrument 2: AGENTS.md Constraint Compliance Evaluation}. A structured test assessing Codex's fidelity to project-level architectural constraints defined in \texttt{AGENTS.md}. The evaluation measured layer boundary adherence across the full Clean Architecture stack, naming convention consistency, build command compliance, and constraint persistence across multi-session and post-compaction execution windows. This instrument directly leverages Codex's fine-tuned instruction-following capability and has no exact analogue in tools lacking a comparable deterministic project instruction mechanism.

\textbf{Instrument 3: External Benchmark Triangulation}. Observed behavior was contextualized against published benchmarks (SWE-Bench Pro, OSWorld-Verified, Terminal-Bench 2.0, Toolathlon, BrowseComp, MMMU Pro), practitioner community sentiment analysis, and comparative tool analyses. All benchmark figures were independently verified against primary sources including BenchLM \cite{benchlm2026}, Vellum AI \cite{vellum2026}, official OpenAI and Anthropic documentation, and the MorphLLM SWE-Bench Pro leaderboard \cite{morphllm2026}.

\textbf{Measurement note}. Rubric scores are practitioner observational ratings from this evaluation period and should be interpreted as internal comparative signals. Benchmark figures sourced from third-party evaluators reflect specific evaluation configurations and scaffolding conditions that may differ across assessments.

\section{Results}

\subsection{Quantitative Rubric Scores}
Table~\ref{tab:category_scores} presents category-level scores; the full 26-criterion breakdown is in Table~\ref{tab:full_rubric}. The weighted mean was 4.50 / 5.00 (90.0\%).

\begin{table}[t]
\caption{Category-Level Rubric Scores (1 to 5 scale)}
\label{tab:category_scores}
\centering
\setlength{\tabcolsep}{4pt}
\begin{tabular}{@{}lr@{}}
\toprule
\textbf{Category} & \textbf{Mean} \\
\midrule
Workflow and Setup & 5.00 \\
Code Generation Quality & 4.50 \\
Context and Memory & 4.40 \\
Architecture and Planning & 4.75 \\
Problem Solving & 5.00 \\
Documentation & 5.00 \\
Operational / Non-Functional & 3.60 \\
Overall Productivity Impact & 5.00 \\
\midrule
\textbf{Overall Weighted Mean} & \textbf{4.50} \\
\bottomrule
\end{tabular}
\end{table}

\begin{table}[t]
\caption{Full 26-Criterion Rubric Scores}
\label{tab:full_rubric}
\centering
\setlength{\tabcolsep}{3pt}
\begin{tabular}{@{}llc@{}}
\toprule
\textbf{Category} & \textbf{Criterion} & \textbf{Score} \\
\midrule
\multirow{2}{*}{Workflow \& Setup}
 & Setup \& Configuration & 5 \\
 & Learning Curve & 5 \\
\cmidrule(l){2-3}
\multirow{4}{*}{Code Generation}
 & Accuracy \& Correctness & 5 \\
 & Feature Creation & 5 \\
 & Maintainability & 4 \\
 & Code Consistency & 4 \\
\cmidrule(l){2-3}
\multirow{5}{*}{Context \& Memory}
 & Instruction Following & 4 \\
 & Context Awareness & 5 \\
 & Context Window Limits & 4 \\
 & Token Efficiency & 4 \\
 & Project-Level Instructions & 5 \\
\cmidrule(l){2-3}
\multirow{4}{*}{Arch.\ \& Planning}
 & Feature Planning & 4 \\
 & Architecture Summary & 5 \\
 & Multi-Step / Role-Based & 5 \\
 & Sub-Agent Support & 5 \\
\cmidrule(l){2-3}
\multirow{4}{*}{Problem Solving}
 & Bug Finding \& Fixing & 5 \\
 & Iterative Refinement & 5 \\
 & Handling AI Mistakes & 5 \\
 & Tech Debt Identification & 5 \\
\cmidrule(l){2-3}
\multirow{5}{*}{Operational}
 & Security \& IP & 1 \\
 & Cost Considerations & 4 \\
 & Response Time / Latency & 3 \\
 & IDE Integration & 5 \\
 & Rate Limiting & 5 \\
\cmidrule(l){2-3}
Documentation & Documentation \& Comments & 5 \\
\cmidrule(l){2-3}
Overall Impact & Productivity Observations & 5 \\
\bottomrule
\end{tabular}
\end{table}

\subsection{External Benchmark Triangulation}
Table~\ref{tab:benchmarks} presents Codex's performance across major AI coding and agentic benchmarks, compared against Claude Opus 4.6. All figures are sourced from independent evaluators and cross-verified against official primary documentation.

\begin{table}[t]
\caption{Codex (GPT-5.4) vs.\ Claude Opus 4.6: Key Benchmark Results}
\label{tab:benchmarks}
\centering
\setlength{\tabcolsep}{3pt}
\begin{tabular}{@{}lrrl@{}}
\toprule
\textbf{Benchmark} & \textbf{Codex} & \textbf{Opus 4.6} & \textbf{Domain} \\
\midrule
SWE-Bench Verified & $\sim$80.0\% & 81.4\% & Standard issue resolution \\
SWE-Bench Pro & \textbf{57.7\%} & 53.4\% & Novel engineering logic \\
OSWorld-Verified & \textbf{75.0\%} & 72.7\% & Desktop automation (Human: 72.4\%) \\
Terminal-Bench 2.0 & \textbf{75.1\%} & 65.4\% & CLI \& agentic terminal ops \\
Toolathlon & \textbf{54.6\%} & 47.2\% & External tool invocation \\
BrowseComp & 82.7\% & \textbf{84.0\%} & Multi-step agentic web search \\
MMMU Pro (w/ tools) & \textbf{81.2\%} & 77.3\% & Visual \& multimodal reasoning \\
MRCR v2 (1M context) & $\sim$70.0\%$^\dagger$ & \textbf{76.0\%} & Long-context retrieval \\
\bottomrule
\multicolumn{4}{@{}l}{\footnotesize $^\dagger$GPT-5.4 1M-context MRCR score unverifiable from indexed sources; approximate.}
\end{tabular}
\end{table}

\subsection{AGENTS.md Constraint Compliance Evaluation}
Codex demonstrated strong and consistent fidelity to project-level architectural constraints specified in \texttt{AGENTS.md}. The deterministic instruction hierarchy---CLI flags $\rightarrow$ project-scoped \texttt{AGENTS.md} files $\rightarrow$ global user configuration $\rightarrow$ system defaults---was respected across all evaluation sessions, including post-compaction segments where free-form session context had been pruned.

Key findings from the compliance evaluation:
\begin{itemize}
    \item Clean Architecture layer boundaries were enforced without per-prompt reminders across the 24-file feature implementation
    \item \texttt{AGENTS.md}-specified naming conventions and state management patterns persisted through context compaction events; file-based instructions proved significantly more durable than equivalent free-form session instructions
    \item Negative constraints (``do not modify files outside this directory'') were honored in Worktree and Cloud execution modes, with the git worktree boundary providing an additional technical enforcement layer
    \item Feature Planning scored 4/5 rather than 5/5, reflecting cases where Codex bypassed the planning phase on ambiguous sub-tasks and proceeded directly to code generation---a manifestation of the rushing behavior documented in the practitioner literature \cite{wiegold2026}
\end{itemize}

The enterprise-grade configurability of \texttt{AGENTS.md} hierarchies---allowing global organizational standards at the repository root and hyper-specific overrides at the microservice directory level---was consistently cited by evaluators as a significant operational advantage over the re-prompting conventions required by tools lacking comparable instruction infrastructure.

\subsection{Observed Workflow Behavior}

\subsubsection{Cloud Worktree Execution and Parallel Isolation}
Codex's cloud-native execution model represents the most architecturally distinctive feature among frontier coding tools. Three execution modes are supported: Local (runs on developer machine within the project), Worktree (creates an isolated git worktree on the developer machine), and Cloud (runs in isolated OpenAI-managed containers). The Cloud mode's two-phase execution architecture provides a documented security boundary:
\begin{itemize}
    \item \textbf{Setup Phase}: Internet access enabled for dependency installation; all configured secrets available
    \item \textbf{Agent Phase}: Internet access disabled by default; all secrets removed before execution begins; network proxy enforced for abuse prevention
\end{itemize}

Container state is cached for up to 12 hours and automatically invalidated when setup scripts, environment variables, or secrets change, providing reproducible execution environments without per-task infrastructure overhead. Subagents in Codex run in their own isolated containers with scoped file access and no network access, providing a strict security perimeter for parallel subtask execution.

In the parallel isolation evaluation, five concurrent tasks were initiated across two repositories. Each task operated on an independent git worktree without conflict, with results reviewable as they completed. This operational model---\textit{delegate and review}---has no direct analogue in single-context tools such as Copilot or Cursor.

\subsubsection{Context Compaction and Long-Horizon Execution}
GPT-5.4's context compaction is an architectural capability purpose-trained into the model, not an external summarization layer applied post-hoc. When an execution session approaches the context window limit, the model itself prunes earlier history while preserving task-critical state---architectural decisions, constraint specifications, and implementation commitments---enabling coherent continuation across what would otherwise be session boundaries \cite{openai_compaction}.

In practice, this delivers:
\begin{itemize}
    \item Approximately 20--40\% reduction in total token consumption for sessions exceeding 100,000 tokens
    \item Approximately 30\% fewer thinking tokens for equivalent reasoning effort compared to predecessor models
    \item Documented internal OpenAI runs exceeding 24 hours of continuous autonomous execution
\end{itemize}

The Token Efficiency criterion scored 4/5, reflecting this advantage tempered by user-reported cases where compaction mid-task caused loss of nuanced working context that had not yet been committed to the task plan. Some practitioners describe this as the agent ``forgetting where it was'' in multi-step refactoring sequences, requiring developer re-contextualization \cite{openai_codex_docs}.

The complementary \textbf{Tool Search} feature addresses a different token efficiency problem. Rather than injecting monolithic 4,000-token tool manifests into every API request, GPT-5.4 dynamically loads tool definitions, MCP server schemas, and API manifests on demand during execution. This architectural optimization produces approximately 47\% fewer input tokens for agentic tasks, substantially reducing the Total Cost of Ownership for high-volume automated pipelines \cite{openai_compaction}.

\subsubsection{Fast Mode Velocity}
Codex's Fast Mode delivers up to 1.5$\times$ faster token velocity using the identical underlying model---a genuine throughput improvement rather than a quality-speed trade-off. This feature was consistently cited across independent practitioner reviews as the most meaningful developer experience improvement in the GPT-5.4 release: ``I would usually rather finish the task and spend the quota than save quota while an agent crawls through the job'' \cite{holter2026}.

Fast Mode is most effective for well-defined, operationally clear tasks where exhaustive upfront planning would delay rather than improve outcomes. The Response Time / Latency rubric score of 3/5 reflects the counterpart: on complex multi-file tasks requiring deep reasoning, Codex does not invest the same upfront ``thinking budget'' that Claude Code's architecture favors, occasionally producing implementations requiring more downstream correction cycles. The practical guidance is to reserve Fast Mode for execution phases on well-specified tasks and allow standard mode for planning phases on architectural decisions.

\subsubsection{Analytical Problem Solving and Bug Identification}
A consistent pattern across practitioner evaluations is Codex's exceptional analytical depth when deployed as an autonomous code reviewer or debugger. In direct side-by-side comparisons, practitioners report that for an identical piece of complex software, Claude Code identifies one to two high-level architectural issues while GPT-5.4 excavates four to five granular problems---flagging subtle edge cases, variable type mismatches, implicit nullability violations, and minor state inconsistencies that human reviewers systematically miss \cite{chandler2026}.

This depth has been validated in production environments: Codex has been deployed to debug complex WebGL terminal rendering artifacts and deep OS-level Windows memory crashes with a documented success rate exceeding Claude Opus 4.6 on equivalent tasks \cite{chandler2026}. In the evaluation period, Bug Finding, Iterative Refinement, Handling AI Mistakes, and Tech Debt Identification all achieved the maximum score of 5/5.

Sustained execution capability further amplifies this analytical depth. With a well-specified \texttt{PLANS.md} execution file, Codex has been documented running continuously for over 46 minutes---autonomously cycling through reading files, committing changes, testing logic, debugging its own errors, and pushing code without losing operational context \cite{chandler2026}. For legacy remediation workloads, the model has successfully refactored upwards of 12,000 lines of undocumented Python code in single extended sessions, operating with the methodical patience of a senior engineer and silently resolving latent race conditions that had existed in production environments for months.

\subsubsection{Destructive Execution Risk}
The most operationally significant safety concern identified in practitioner community analysis is what this evaluation terms \textit{destructive execution}: Codex's documented tendency, when given underspecified or contradictory instructions, to rush toward completion rather than pausing to seek clarification \cite{wiegold2026}. In documented cases, Codex has completed a portion of a requested refactoring correctly before overwriting working progress with hallucinatory code, resulting in zero net progress and requiring full session rollback \cite{reddit2026}.

This behavior stands in direct contrast to Claude Code's documented tendency to collaboratively discuss ambiguous intent before executing. The root cause is not a lack of capability but a \textit{failure of evaluation}: Codex does not automatically assess whether its plan is complete before executing it. The full mitigation stack is:
\begin{itemize}
    \item Comprehensive \texttt{AGENTS.md} negative constraints (explicit ``do not'' rules and permission boundaries)
    \item Sandbox mode configured to \texttt{workspace-write} rather than \texttt{danger-full-access} for critical sessions
    \item Approval Policy set to \texttt{on-request} for sessions involving irreversible operations
    \item Git worktree isolation: even if destructive changes are executed within a worktree, the main branch remains unaffected and the worktree can be discarded cleanly
\end{itemize}

The git worktree architecture is particularly valuable as a mitigation layer: it converts a potentially catastrophic irreversible action into a discardable branch operation.

\subsubsection{Rate Limiting Behavior}
In a notable contrast to Claude Code (which scored 1/5 on Rate Limiting in this evaluation series), Codex achieved 5/5, reflecting substantially more generous utilization allowances on standard Plus subscription tiers. For development teams operating on Plus plans, rate limiting was not a meaningful friction point during the evaluation period.

However, the higher-capability Pro tier presents a documented limitation: approximately 40 minutes of reasoning capacity per 5-hour rolling window. A single complex reasoning-intensive task of 20 minutes can consume half the available Pro allowance, making the tier economically challenging for sustained professional workflows \cite{openai_community2026}. The absence of a mid-tier subscription option between the \$20 Plus and \$200 Pro tiers---in contrast to Anthropic's \$100 Max tier---leaves Pro-tier users with limited options when quota is exhausted other than waiting for quota restoration or downgrading to the medium-effort tier.

\subsection{Cross-Model Audit Workflow}
An emergent best practice documented across the practitioner community is the use of Codex as an auditor for Claude-generated code, and vice versa \cite{chandler2026}. GPT-5.4's granular analytical depth makes it exceptionally effective at catching Claude Opus 4.6's occasional over-engineering tendencies, consistently suggesting cleaner, simplified, and more security-conscious implementations. Claude Opus 4.6, in turn, excels at catching architectural inconsistencies and edge case gaps in error handling within Codex-generated code.

OpenAI has operationalized this workflow through the release of an official Codex plugin for Claude Code, enabling direct cross-provider code review within a single session \cite{mindstudio2026}. Running iterative review cycles---Opus 4.6 critically reviews a Codex plan, Codex reviews the revised Opus implementation---consistently produces superior outcomes to either model operating alone \cite{chandler2026}.

\subsection{Documentation Generation}
Documentation generation achieved the maximum score of 5/5. Codex demonstrates strong capability for generating comprehensive, accurate documentation including inline comments, API doc comments, and README files. The Architecture Summary capability (5/5) reflects GPT-5.4's ability to analyze complex systems including hardware-oriented architectures such as Post Fusion systems involving quantized sub-model deployment on embedded inference hardware, producing structured and technically accurate summaries that span both software and hardware engineering domains.

\section{Discussion}

\subsection{The Analytical Depth Advantage}
The finding most clearly supported by both rubric scores and independent benchmark data is Codex's superior analytical depth for novel problem-solving. The 4.3 percentage point lead on SWE-Bench Pro (57.7\% vs.\ 53.4\%) is not a marginal measurement difference---it reflects a qualitatively distinct ability to reason from first principles on problems outside the training distribution, rather than retrieving memorized solutions. In practitioner terms, this means Codex systematically surfaces more bugs per review cycle, identifies a broader range of technical debt, and produces solutions less dependent on analogous patterns in its training data \cite{chandler2026}.

For engineering teams whose primary challenge is understanding and improving large, undocumented legacy codebases---particularly those with complex concurrency logic, deep OS-level interactions, or intricate backend API integrations---this analytical edge represents a compelling production argument for Codex as the execution and debugging tool of choice.

\subsection{Cloud Architecture as a Competitive Differentiator}
Codex's cloud-native execution model, built around git worktree isolation and containerized task execution, provides capabilities with no direct analogue in locally-executed tools. Three advantages stand out for enterprise and team-scale adoption:

\textbf{True parallel project execution}: Multiple independent tasks across multiple repositories run simultaneously, each in an isolated container, without contention for local compute resources. A developer can initiate parallel refactoring tasks across distinct microservices and review results as they complete---an operational model that single-context tools cannot support.

\textbf{Reproducible execution environments}: Container state is deterministically constructed from setup scripts and cached, providing reproducible execution that local filesystem-dependent tools cannot match. This is directly relevant for regulated environments requiring audit trails and reproducible builds.

\textbf{Documented security boundary}: The two-phase execution model (internet-enabled setup, offline agent execution) provides a clear security perimeter. Code exfiltration during the agent execution phase is architecturally prevented, addressing a portion of the IP concern that drives Security \& IP to score 1/5 under general evaluation conditions.

\subsection{Native Compaction vs.\ Session Architecture Discipline}
Competing tools manage long-context sessions through architectural workarounds: Claude Code uses sub-agent context isolation and developer-initiated summarization prompts. Codex's native compaction takes a fundamentally different approach: the model itself manages its own context window, pruning history while preserving task-critical state without developer intervention.

The practical consequence is that Codex handles context window transitions more transparently---developers do not need to architect sessions around context limits or manually instruct summarization. The 47\% token reduction from Tool Search compounds this advantage, making Codex the cost-optimal choice for high-volume automated pipelines where total token consumption is the primary economic variable.

The primary operational caveat remains: compaction can occasionally discard nuanced working context not yet committed to the task plan. Teams running complex refactoring sessions should use \texttt{PLANS.md} execution files to externalize task state explicitly, providing a durable record that survives context compaction events.

\subsection{The Destructive Execution Problem and Its Mitigations}
The \textit{destructive execution} pattern---Codex's tendency to rush forward with assumptions rather than seeking clarification---is a real and documented friction point \cite{wiegold2026,reddit2026}. It is important to characterize this accurately: this is not a hallucination problem in the traditional sense but a \textit{failure of evaluation}. Codex demonstrably understands the code it is working with; it simply does not automatically assess whether its plan is fully specified before executing it.

The mitigation is clear and well-defined. The combination of comprehensive \texttt{AGENTS.md} authoring, \texttt{workspace-write} sandbox mode, \texttt{on-request} approval policy, and git worktree isolation creates a defense-in-depth architecture that contains the impact of any individual destructive execution event to a discardable branch. Organizations that invest in this mitigation stack report a substantially improved experience relative to those using Codex in default configuration on underspecified tasks.

\subsection{Economic Considerations and Total Cost of Ownership}
On raw API pricing, Codex Standard is substantially more cost-effective than Claude Opus 4.6: 50\% cheaper on input tokens (\$2.50 vs.\ \$5.00 per million) and 40\% cheaper on output tokens (\$15.00 vs.\ \$25.00 per million). At the cached input tier, Codex (\$0.25/M) is half the cost of Opus 4.6 (\$0.50/M), making iterative development on large codebases significantly more economical.

When combined with the 47\% token reduction from Tool Search, Codex Standard represents the cost-optimal choice for high-volume automated testing, proof-of-concept generation, and bulk refactoring workloads with well-specified \texttt{AGENTS.md} guardrails. Table~\ref{tab:pricing} presents the current pricing landscape.

\begin{table}[t]
\caption{API Pricing Comparison (per 1M tokens, verified April 2026)}
\label{tab:pricing}
\centering
\setlength{\tabcolsep}{4pt}
\begin{tabular}{@{}lrrrr@{}}
\toprule
\textbf{Model} & \textbf{Input} & \textbf{Output} & \textbf{Cached} & \textbf{Context} \\
\midrule
Codex (GPT-5.4) & \$2.50 & \$15.00 & \$0.25 & 1.05M \\
Claude Opus 4.6 & \$5.00 & \$25.00 & \$0.50 & 1.00M \\
Codex Pro & \$30.00 & \$180.00 & N/A & $>$272K \\
\bottomrule
\end{tabular}
\end{table}

However, Total Cost of Ownership (TCO) shifts considerably at scale. Codex's lower accuracy on complex multi-file coordinated tasks---evidenced by its $\sim$1.4 percentage point deficit on SWE-Bench Verified relative to Claude Opus 4.6---leads to higher rates of automated retries, destructive execution rollbacks, and required developer debugging intervention. When the senior developer hourly cost of untangling a hallucinatory implementation is factored in, TCO can favor the more expensive but more architecturally reliable Claude Opus 4.6 for novel multi-file work where architectural integrity is the primary constraint.

\subsection{The Dual-Wielding Recommendation}
The most well-corroborated finding from practitioner community analysis is that the optimal deployment strategy is not a single-model selection but a \textit{dual-wielding} workflow \cite{chandler2026}. Claude Opus 4.6 and Codex possess highly asymmetric capability profiles that render them strategic complements rather than substitutes:

\begin{itemize}
    \item \textbf{Claude Opus 4.6 as architect}: UI generation, complex multi-file architectural refactoring, ambiguous task handling, proactive QA dispatch via Agent Teams parallel orchestration
    \item \textbf{Codex as executor and analytical auditor}: Backend logic implementation, concurrency debugging, technical debt excavation, cross-model code review of Claude-generated code, and high-volume automated pipeline execution
\end{itemize}

Engineering organizations operating at the frontier of development efficiency route tasks dynamically based on their cognitive requirements: architectural depth and creativity to Claude, execution velocity and analytical verification to Codex. This dual-wielding workflow is directly enabled by OpenAI's official Codex plugin for Claude Code, allowing cross-provider review within a single integrated session \cite{mindstudio2026}.

\subsection{Terminal Regression: A Practical Caveat}
An operationally relevant finding for DevOps-focused teams is that GPT-5.4 scores 75.1\% on Terminal-Bench 2.0---lower than its specialized predecessor GPT-5.3-Codex at 77.3\% \cite{nxcode2026}. The unification of GPT-5.4 with the Codex model provided substantial gains in visual reasoning and multimodal capability but introduced a measurable regression in raw terminal fidelity. Teams with workflows heavily dependent on SSH operations, complex Git automation, shell script debugging, and build system manipulation should note that GPT-5.3-Codex remains the superior option for pure terminal-centric workloads. For mixed workflows combining terminal operations with code generation, the generalization advantages of GPT-5.4 outweigh the terminal regression.

\section{Limitations}

\textbf{Single project, limited task scope}. All conclusions are drawn from an evaluation in one codebase. Generalizability to different languages, frameworks, or team sizes is not established.

\textbf{No controlled comparison}. Without a direct controlled comparison against competing tools on identical tasks, relative productivity claims cannot be causally attributed to Codex specifically.

\textbf{Benchmark sourcing and scaffolding variability}. Several benchmarks, particularly SWE-Bench Pro and GDPval, are reported from secondary aggregator sources where the exact evaluation scaffolding may differ from primary published conditions. GDPval figures are reported as percentages in secondary sources but the primary scoring system is Elo-based (GPT-5.4: 1672; Claude Opus 4.6: 1606); percentage representations are approximations. The MRCR v2 score at 1M context for GPT-5.4 could not be verified from any indexed source at time of writing and should be treated as an approximation.

\textbf{Snapshot validity}. Codex and its cloud execution infrastructure are actively evolving. Rate limiting, pricing, compaction behavior, and benchmark positions may change with future releases. Terminal-Bench regression may be addressed in future GPT-5.4 updates.

\textbf{Destructive execution variability}. The destructive execution behavior is prompt- and task-dependent. Its frequency in production contexts will vary substantially based on \texttt{AGENTS.md} quality, sandbox configuration, and task specificity. Evaluations conducted under tight constraint specifications will observe significantly lower incidence than those conducted on loosely defined tasks.

\section{Conclusion}
This paper evaluated Codex (GPT-5.4) as an AI coding assistant through a practitioner benchmark grounded in production feature implementation. We address our three research questions as follows:

\textbf{Q1 (Core Capability)}. Codex functions effectively as an AI coding assistant for structured, convention-driven codebases, with particular strength in analytical problem solving, technical debt identification, and sustained execution. Its SWE-Bench Pro score of 57.7\% reflects a genuine and verified superiority in novel engineering reasoning. The 24-file feature implementation was completed with full Clean Architecture integrity, and Bug Finding, Iterative Refinement, Handling AI Mistakes, and Tech Debt Identification all achieved the maximum score of 5/5. The primary caveat is the destructive execution tendency under underspecified prompts, which requires investment in \texttt{AGENTS.md} authoring and sandbox configuration to mitigate effectively.

\textbf{Q2 (Operational Characteristics)}. Codex's operational profile is defined by three distinctive characteristics: cloud-native sandboxed execution via git worktrees (enabling parallel isolation without consuming local developer resources), purpose-trained context compaction (enabling 24+ hour sessions with $\sim$30\% fewer thinking tokens and $\sim$47\% reduced input token consumption via Tool Search), and Fast Mode velocity (up to 1.5$\times$ token throughput for well-defined tasks). Rate limiting is not a friction point on standard Plus tiers (5/5), a significant advantage over competing tools. The primary operational gaps are the Pro tier reasoning quota ceiling and the Security \& IP score (1/5), which reflects the inherent constraints of transmitting proprietary code to cloud-hosted model infrastructure.

\textbf{Q3 (Architectural Differentiators)}. Codex's git worktree architecture, cloud execution model, and native compaction collectively provide competitive capabilities with no direct analogues in locally-executed tools. True parallel task execution across multiple repositories, reproducible containerized environments, and the Tool Search token reduction position Codex as the cost-optimal choice for high-volume automated pipelines. The official Codex plugin for Claude Code enables the dual-wielding workflow that practitioners consistently identify as the highest-productivity deployment pattern, routing architectural creativity to Claude Opus 4.6 and execution depth to Codex.

The overall finding is that Codex (GPT-5.4) is the apex execution engine and analytical debugger in the current frontier coding tool landscape. Its strengths---analytical depth on novel problems, cloud-native parallelism, token efficiency at scale, and first-in-class desktop automation---are precisely the capabilities that complement Claude Code's architectural creativity and parallel orchestration strengths. For teams willing to invest in the operational discipline of \texttt{AGENTS.md} authoring and sandbox configuration, Codex represents a compelling productivity multiplier whose analytical capabilities, economic efficiency, and parallel execution architecture make it an essential component of any serious agentic development workflow.

\begin{thebibliography}{00}
\setlength{\itemsep}{-2pt}

\bibitem{chen2021evaluating}
M.~Chen et~al., ``Evaluating Large Language Models Trained on Code,'' arXiv preprint arXiv:2107.03374, 2021.

\bibitem{jimenez2024swebench}
C.~E.~Jimenez et~al., ``SWE-bench: Can Language Models Resolve Real-World GitHub Issues?'' Proc.\ International Conference on Learning Representations (ICLR), 2024.

\bibitem{fryer2024hallucination}
D.~Fryer et~al., ``Do Large Language Models Know What They Don't Know?'' arXiv preprint arXiv:2402.10768, 2024.

\bibitem{peng2023copilot}
S.~Peng, E.~Kalliamvakou, P.~Cihon, and M.~Demirer, ``The Impact of AI on Developer Productivity: Evidence from GitHub Copilot,'' arXiv preprint arXiv:2302.06590, 2023.

\bibitem{mckinsey2023}
McKinsey Global Institute, ``The Economic Potential of Generative AI: The Next Productivity Frontier,'' McKinsey and Company, Jun.\ 2023.

\bibitem{metr2025}
METR, ``Measuring the Impact of Early-2025 AI on Experienced Open-Source Developer Productivity,'' METR Technical Report, 2025.

\bibitem{creswell2018mixed}
J.~W.~Creswell and V.~L.~Plano Clark, \textit{Designing and Conducting Mixed Methods Research}, 3rd~ed.\ Thousand Oaks, CA: SAGE Publications, 2018.

\bibitem{morphllm2026}
MorphLLM, ``SWE-Bench Pro Leaderboard,'' 2026. [Online]. Available: https://www.morphllm.com/swe-bench-pro

\bibitem{nxcode2026}
NxCode, ``GPT-5.4 Complete Guide 2026: Features, Pricing, Benchmarks \& How to Use,'' 2026. [Online]. Available: https://www.nxcode.io/resources/news/gpt-5-4-complete-guide-features-pricing-models-2026

\bibitem{benchlm2026}
BenchLM, ``GPT-5.4 Benchmark Results,'' 2026. [Online]. Available: https://benchlm.ai/models/gpt-5-4

\bibitem{vellum2026}
Vellum AI, ``Claude Opus 4.6 Benchmarks and Model Card,'' 2026. [Online]. Available: https://www.vellum.ai/blog/claude-opus-4-6-benchmarks

\bibitem{chandler2026}
C.~Nguyen, ``Codex with GPT-5.4 vs Claude Code with Opus 4.6 --- Why I Now Use Both,'' 2026. [Online]. Available: https://chandlernguyen.com/blog/2026/03/13/codex-gpt-5-4-vs-claude-code-opus-4-6-dual-wielding-ai-coding-tools/

\bibitem{holter2026}
A.~Holter, ``GPT-5.4 Fast Mode Is the Best Part of the Release,'' 2026. [Online]. Available: https://adam.holter.com/gpt-5-4-fast-mode-is-the-best-part-of-the-release/

\bibitem{wiegold2026}
T.~Wiegold, ``I Tested GPT 5.4 Against Every Rival --- Here's My Honest Review,'' 2026. [Online]. Available: https://thomas-wiegold.com/blog/i-tested-gpt-5-4-against-every-rival/

\bibitem{reddit2026}
OpenAI Developer Community, ``GPT-5.4 in Codex Feels Degraded,'' 2026. [Online]. Available: https://community.openai.com/t/gpt-5-4-in-codex-feels-degraded/1378747

\bibitem{openai_community2026}
OpenAI Developer Community, ``Understanding the New Codex Limit System After the April 9 Update,'' 2026. [Online]. Available: https://community.openai.com/t/understanding-the-new-codex-limit-system-after-the-april-9-update/1378768

\bibitem{mindstudio2026}
MindStudio, ``OpenAI Codex Plugin for Claude Code: Cross-Provider Review,'' 2026. [Online]. Available: https://www.mindstudio.ai/blog/openai-codex-plugin-claude-code-cross-provider-review

\bibitem{openai_codex_docs}
OpenAI, ``Codex Developer Documentation,'' 2026. [Online]. Available: https://developers.openai.com/codex

\bibitem{openai_cloud_envs}
OpenAI, ``Cloud Environments --- Codex Web,'' 2026. [Online]. Available: https://developers.openai.com/codex/cloud/environments

\bibitem{openai_compaction}
OpenAI, ``Building More with GPT-5.1-Codex-Max,'' 2026. [Online]. Available: https://openai.com/index/gpt-5-1-codex-max/

\end{thebibliography}

\end{document}
